%%
%% techbrief-template.tex -- a blank ACM TechBrief, ready to fill in.
%%
%% Copy this file (and techbrief-template.bib) under a new name, replace
%% every ALL-CAPS placeholder, and build with
%%
%%     make DOC=<your-file-name-without-.tex>
%%
%% Everything marked "CHANGE" is per issue; everything else is part of
%% the TechBrief design and normally left alone.  README.md documents
%% every command; techbrief-sample.tex is a complete worked example
%% (Issue 18).
%%
%% Uses the unmodified acmart class plus techbrief.sty.  Compile with
%% LuaLaTeX and biber (the Makefile does this for you).
%%
\documentclass[acmlarge,nonacm,natbib=false]{acmart}
\usepackage{techbrief}

%% CHANGE: your bibliography file(s), one \addbibresource per file.
\addbibresource{techbrief-template.bib}

%% CHANGE: DOI and ISBN, printed at the foot of the back page.
\acmDOI{10.1145/XXXXXXX}
\acmISBN{979-8-4007-XXXX-X}

%% CHANGE: season and issue number, printed in every page footer.
\techbriefissue{SEASON YEAR}{NN}

%% CHANGE: acknowledgments, printed in the panel on the back page.
\techbriefacknowledgments{The authors would like to thank REVIEWERS for
  their reviews and input. A thank you to the ACM TechBrief Committee and
  the ACM Technology Policy Council for their valuable input.}

%% OPTIONAL: cover art for this issue (202.9bp wide x 201.1bp high).
%% Without this line the Issue 18 cover is used.  Prefer vector art: a
%% TikZ picture, or a PDF via \includegraphics.
% \techbriefcover{\includegraphics[width=202.9bp,height=201.1bp]{cover.pdf}}

\begin{document}

%% CHANGE: the title.  The optional argument is the short title used in
%% the running head on pages 2 onward (keep it under about 40
%% characters).  Use \\ to choose where the two-line banner title breaks.
\title[SHORT TITLE FOR THE RUNNING HEAD]{FULL TITLE OF THE TECHBRIEF,\\
  ON NO MORE THAN TWO LINES}

%% CHANGE: authors, one \author per person, in byline order.  The byline
%% reads "by A, B, and C".
\author{FIRST AUTHOR}
\author{SECOND AUTHOR}
\author{THIRD AUTHOR}

\maketechbrief

\begin{techbriefbody}
%% The opening paragraphs: what the technology is and why it matters to
%% policy makers.  Cite sources with \cite{key}; citations print as
%% superscripts and are numbered in the order they are first cited.
INTRODUCTORY PARAGRAPH. State the issue in plain language and define
any technical terms.\cite{example2025report}

SECOND PARAGRAPH. Explain why the issue matters now.

%% A table (optional).  Column widths must add up to 515.6bp; each L{}
%% column includes a 6bp gutter on its right.  End the header row with
%% \tbheadend and every other row with \\.
\begin{techbrieftable}{TABLE CAPTION}
    {L{115.3bp}L{171.7bp}L{228.6bp}}
  \rowcolor{tbblue}
    & \tbhead{COLUMN A} & \tbhead{COLUMN B} \tbheadend
  \tbrowhead{ROW ONE}   & CELL & CELL \\
  \tbrowhead{ROW TWO}   & CELL & CELL \\
\end{techbrieftable}

\section*{Key Policy Dimensions}

\runinhead{FIRST DIMENSION.} PARAGRAPH introducing the first policy
dimension.

\begin{itemize}
\item FIRST POINT.\cite{example2024article}
\item SECOND POINT.
  \begin{itemize}
  \item A SUB-POINT (one level of nesting is supported).
  \end{itemize}
\end{itemize}

%% An italic summary sentence.  Use \begin{takeaway}[12bp] to indent it
%% to the bullet text.
\begin{takeaway}
ONE-SENTENCE TAKEAWAY for this dimension.
\end{takeaway}

\runinhead{SECOND DIMENSION.} PARAGRAPH.

\section*{Conclusion}

CONCLUDING PARAGRAPHS summarizing the policy questions raised.
\end{techbriefbody}

%%
%% The back page: references at the top, and the Additional Information /
%% Acknowledgments panel at the foot (filled in automatically).
\techbriefbackpage
\printbibliography

\end{document}
\endinput
%%
%% End of file `techbrief-template.tex'.
